\section{提示词：把愿望写成可以验收的任务}
\label{sec:ai-prompts}

刚才请 AI 修改计分器时，我们给了原始输出、新规则和允许改动的范围。换一个问题，你可能会问：“我说得挺清楚，它怎么还是改错了？”先看看你的提示词有没有把\textbf{现象、目标与验收}分开。“优化识别效果”只有方向，没有说明在哪组图上、用什么指标判断、允许付出多少延迟；模型只好替你猜。猜得再流畅，也不等于猜中了队内需求。

给 AI 的请求和材料通常叫提示词（Prompt）。对于开发任务，先用下面几项把问题说清楚，让 AI 拥有充足的上下文。
\begin{center}
\begin{tabular}{p{0.16\textwidth}p{0.73\textwidth}}
\toprule
信息 & 应写出的内容 \\
\midrule
目标与现象 & 现在发生什么，期望变成什么；把“修复卡顿”变成一个可观察行为。\\
相关证据 & 仓库、入口文件、复现命令、完整的相关报错、输入样本与版本。\\
范围与约束 & 哪个模块可修改，接口、单位、依赖与硬件条件有什么限制。\\
验收与交付 & 用什么命令、样本和指标判断完成；交付代码、说明还是方案。\\
不确定处 & 哪些条件还不知道，先调查什么，哪些决策需要向你确认。\\
\bottomrule
\end{tabular}
\end{center}

但也不要为了凑“上下文”粘贴几万行日志。给出首次错误附近的日志、运行命令及文件路径，让有工具的代理自行定位其余内容；未获本地文件访问能力的聊天需要你附上最小相关代码或文件。截图适合表达界面、曲线与视觉现象，编译错误优先附可复制的文本。发送前检查日志与附件，移除令牌、密码、队友隐私和无关的队内资料。

\begin{example}[同一个错误，怎样从许愿变成任务]
假设你在计分器里把 bonus 改成了 -5，运行时却仍看到 75。与其只说“帮我修一下”，可以先这样请 AI 帮你排查：
\begin{quote}
我在计分器的 \texttt{main.cpp} 中把 bonus 从 15 改成了 -5。基础分是 60，预期输出 55，但点运行仍显示 75。我还没确认这次启动的是不是新构建的目标。请先指出最值得核对的两步，并说明每步要看什么。先给检查办法，保留计分函数。
\end{quote}
这个提问给出了预期和观察结果，也说清了尚未确认的条件。你可以先看构建日志是否显示编译成功，再检查 IDE 实际启动的目标与构建目录。若仍找不到原因，补上相关日志和路径，把调查缩小到能复现的那一步。拿到建议后，用实际结果排除原因。
\end{example}

反馈也要有证据。不要反复发送“还是不行”“再聪明一点”；写清“新构建后输入为……，实际输出……，预期……，这是命令和日志”。如果你怀疑它的方案，要求它比较候选原因、列出能区分这些原因的实验，或指出结论对应的文件与官方资料。\textbf{要让 AI 给出可检查的解释与依据}，不能把 AI 的幻觉当成事实，更不能把一大段解释当成验证。

代理运行时，发现关键约束漏了就及时补充；下一项独立任务可以排队。不要一边让它修函数，一边在同一聊天中突然讨论装系统。若目标已经变了，要明确说“停止当前任务，保存现状，改做……”，并检查仍在运行的命令。功能入口见第~\ref{sec:ai-workflow}~节。

\begin{exercise}\label{ex:ai-prompt}
用自己遇到的一个构建或运行问题练习。如果手边刚好没有错误，可在计分器练习副本删掉一处分号，重新构建并记录第一条相关报错。写一条让 AI 先给检查建议的请求，说明操作步骤、预期、实际现象、相关代码和判断是否恢复的办法。修复后记录实际结果。
\end{exercise}
